\documentclass[10pt,a4paper]{amsart}
\usepackage{amsmath,amssymb,amsthm,graphicx}
\usepackage{verbatim}
\newenvironment{reference}{\comment}{\endcomment}
%\newenvironment{reference}{}{}
\newenvironment{slogan}{\comment}{\endcomment}
\newenvironment{history}{\comment}{\endcomment}

% For commutative diagrams we use Xy-pic
\usepackage[all]{xy}

% We use 2cell for 2-commutative diagrams.
\xyoption{2cell}
% Theorem environments.
%
\theoremstyle{plain}
\newtheorem{theorem}[subsection]{Theorem}
\newtheorem{proposition}[subsection]{Proposition}
\newtheorem{lemma}[subsection]{Lemma}

\theoremstyle{definition}
\newtheorem{definition}[subsection]{Definition}
\newtheorem{example}[subsection]{Example}
\newtheorem{exercise}[subsection]{Exercise}
\newtheorem{situation}[subsection]{Situation}

\theoremstyle{remark}
\newtheorem{remark}[subsection]{Remark}
\newtheorem{remarks}[subsection]{Remarks}

\numberwithin{equation}{subsection}

% Macros
%
\def\lim{\mathop{\mathrm{lim}}\nolimits}
\def\colim{\mathop{\mathrm{colim}}\nolimits}
\def\Spec{\mathop{\mathrm{Spec}}}
\def\Hom{\mathop{\mathrm{Hom}}\nolimits}
\def\Ext{\mathop{\mathrm{Ext}}\nolimits}
\def\SheafHom{\mathop{\mathcal{H}\!\mathit{om}}\nolimits}
\def\SheafExt{\mathop{\mathcal{E}\!\mathit{xt}}\nolimits}
\def\Sch{\mathit{Sch}}
\def\Mor{\mathop{\mathrm{Mor}}\nolimits}
\def\Ob{\mathop{\mathrm{Ob}}\nolimits}
\def\Sh{\mathop{\mathit{Sh}}\nolimits}
\def\NL{\mathop{N\!L}\nolimits}
\def\CH{\mathop{\mathrm{CH}}\nolimits}
\def\proetale{{pro\text{-}\acute{e}tale}}
\def\etale{{\acute{e}tale}}
\def\QCoh{\mathit{QCoh}}
\def\Ker{\mathop{\mathrm{Ker}}}
\def\Im{\mathop{\mathrm{Im}}}
\def\Coker{\mathop{\mathrm{Coker}}}
\def\Coim{\mathop{\mathrm{Coim}}}

% Boxtimes
%
\DeclareMathSymbol{\boxtimes}{\mathbin}{AMSa}{"02}

%
% Macros for moduli stacks/spaces
%
\def\QCohstack{\mathcal{QC}\!\mathit{oh}}
\def\Cohstack{\mathcal{C}\!\mathit{oh}}
\def\Spacesstack{\mathcal{S}\!\mathit{paces}}
\def\Quotfunctor{\mathrm{Quot}}
\def\Hilbfunctor{\mathrm{Hilb}}
\def\Curvesstack{\mathcal{C}\!\mathit{urves}}
\def\Polarizedstack{\mathcal{P}\!\mathit{olarized}}
\def\Complexesstack{\mathcal{C}\!\mathit{omplexes}}
% \Pic is the operator that assigns to X its picard group, usage \Pic(X)
% \Picardstack_{X/B} denotes the Picard stack of X over B
% \Picardfunctor_{X/B} denotes the Picard functor of X over B
\def\Pic{\mathop{\mathrm{Pic}}\nolimits}
\def\Picardstack{\mathcal{P}\!\mathit{ic}}
\def\Picardfunctor{\mathrm{Pic}}
\def\Deformationcategory{\mathcal{D}\!\mathit{ef}}
\usepackage{cleveref}
\crefname{originalsection}{original section}{original sections}
\Crefname{originalsection}{Original section}{Original sections}
\makeatletter
\newlabel{constructions-section-projective-space}{{01ND}{source}{}{sourcepointer-1}{}}
\newlabel{constructions-section-projective-space@cref}{{[originalsection][0][]01ND}{[1][0][]source}{}{}{}}
\makeatother
\crefname{originallemma}{original lemma}{original lemmas}
\Crefname{originallemma}{Original lemma}{Original lemmas}
\makeatletter
\newlabel{descent-lemma-descending-property-closed-immersion}{{02L6}{source}{}{sourcepointer-2}{}}
\newlabel{descent-lemma-descending-property-closed-immersion@cref}{{[originallemma][0][]02L6}{[1][0][]source}{}{}{}}
\makeatother
\crefname{originalreference}{original reference}{original references}
\Crefname{originalreference}{Original reference}{Original references}
\makeatletter
\newlabel{lemma-contracting-rational-tails}{{0E3L}{source}{}{sourcepointer-3}{}}
\newlabel{lemma-contracting-rational-tails@cref}{{[originalreference][0][]0E3L}{[1][0][]source}{}{}{}}
\makeatother
\crefname{originalreference}{original reference}{original references}
\Crefname{originalreference}{Original reference}{Original references}
\makeatletter
\newlabel{lemma-no-rational-tail}{{0E3J}{source}{}{sourcepointer-4}{}}
\newlabel{lemma-no-rational-tail@cref}{{[originalreference][0][]0E3J}{[1][0][]source}{}{}{}}
\makeatother
\crefname{originalreference}{original reference}{original references}
\Crefname{originalreference}{Original reference}{Original references}
\makeatletter
\newlabel{lemma-rational-bridge-zero}{{0E64}{source}{}{sourcepointer-5}{}}
\newlabel{lemma-rational-bridge-zero@cref}{{[originalreference][0][]0E64}{[1][0][]source}{}{}{}}
\makeatother
\crefname{originalreference}{original reference}{original references}
\Crefname{originalreference}{Original reference}{Original references}
\makeatletter
\newlabel{lemma-rational-tail-negative}{{0E63}{source}{}{sourcepointer-6}{}}
\newlabel{lemma-rational-tail-negative@cref}{{[originalreference][0][]0E63}{[1][0][]source}{}{}{}}
\makeatother
\crefname{originalreference}{original reference}{original references}
\Crefname{originalreference}{Original reference}{Original references}
\makeatletter
\newlabel{lemma-rr}{{0BS6}{source}{}{sourcepointer-7}{}}
\newlabel{lemma-rr@cref}{{[originalreference][0][]0BS6}{[1][0][]source}{}{}{}}
\makeatother
\crefname{originalreference}{original reference}{original references}
\Crefname{originalreference}{Original reference}{Original references}
\makeatletter
\newlabel{lemma-vanishing-on-gorenstein}{{0E3E}{source}{}{sourcepointer-8}{}}
\newlabel{lemma-vanishing-on-gorenstein@cref}{{[originalreference][0][]0E3E}{[1][0][]source}{}{}{}}
\makeatother
\crefname{originallemma}{original lemma}{original lemmas}
\Crefname{originallemma}{Original lemma}{Original lemmas}
\makeatletter
\newlabel{varieties-lemma-ampleness-in-terms-of-degrees-components}{{0B5Y}{source}{}{sourcepointer-9}{}}
\newlabel{varieties-lemma-ampleness-in-terms-of-degrees-components@cref}{{[originallemma][0][]0B5Y}{[1][0][]source}{}{}{}}
\makeatother
\crefname{originallemma}{original lemma}{original lemmas}
\Crefname{originallemma}{Original lemma}{Original lemmas}
\makeatletter
\newlabel{varieties-lemma-proper-geometrically-reduced-global-sections}{{0BUG}{source}{}{sourcepointer-10}{}}
\newlabel{varieties-lemma-proper-geometrically-reduced-global-sections@cref}{{[originallemma][0][]0BUG}{[1][0][]source}{}{}{}}
\makeatother
\crefname{originallemma}{original lemma}{original lemmas}
\Crefname{originallemma}{Original lemma}{Original lemmas}
\makeatletter
\newlabel{varieties-lemma-variant-separate-points-tangent-vectors}{{0E8T}{source}{}{sourcepointer-11}{}}
\newlabel{varieties-lemma-variant-separate-points-tangent-vectors@cref}{{[originallemma][0][]0E8T}{[1][0][]source}{}{}{}}
\makeatother
\title{For a nodal curve of genus at least 2 with ample dualizing sheaf, the tricanonical sheaf is very ample and has no first cohomology}
\author{The Stacks Project Authors}
\date{stacks-project:curves.tex:local-snapshot-20261012; GNU FDL 1.2 or later}
\begin{document}
\maketitle
\noindent\textit{Editorial note.} All delivered mathematical text is the author's own, copied from the pinned Stacks Project chapter file curves.tex; no statement or proof step was rewritten, reordered or removed.\par\medskip
\noindent\textit{Editorial note.} The counted proof is the author's complete argument: the global-generation step, the descent to the algebraic closure, the full inequality reduction including both exceptional cases (ampleness excluding degree zero; Riemann-Roch handling Y = X), and the concluding vanishing argument are all retained verbatim.\par\medskip
\noindent\textit{Editorial note.} The label lemma-tricanonical names the counted statement, and the proof environment immediately following it in the source is delivered verbatim as its proof.\par\medskip
\noindent\textit{Editorial note.} The counted statement is carried by the author's own lemma environment in the chapter source.\par\medskip
\noindent\textit{Allowed standard prerequisites (editorial).}\begin{itemize}
\item No rational tails or bridges under ampleness. Varieties, Lemma varieties-lemma-ampleness-in-terms-of-degrees-components; Same chapter, Lemmas lemma-rational-tail-negative and lemma-rational-bridge-zero. Shows X contains no rational tails or bridges, which also rules out degree-zero restrictions of omega\_X to reduced connected subschemes.
\item Global generation of the tricanonical sheaf. Same chapter, Lemma lemma-contracting-rational-tails. Gives that omega\_X\textasciicircum{}tensor 3 is globally generated once tails and bridges are absent.
\item Separation criterion and descent to the algebraic closure. Varieties, Lemma varieties-lemma-variant-separate-points-tangent-vectors; Descent, Lemma descent-lemma-descending-property-closed-immersion. Reduces closed immerseness to surjectivity of H\textasciicircum{}0(X, L) -> H\textasciicircum{}0(Z, L|\_Z) for every degree-2 subscheme Z, and lets one replace k by its algebraic closure.
\item Vanishing on Gorenstein curves and the tail inequality. Same chapter, Lemmas lemma-vanishing-on-gorenstein and lemma-no-rational-tail; Duality for Schemes, Section duality-section-dualizing (Gorenstein hypothesis). Turns the surjectivity problem into the inequality 3 deg(omega\_X|\_Y) > -2 chi(Y, O\_Y) + deg(Z cap Y) and supplies the bound away from the exceptional cases.
\end{itemize}\medskip
\ifdefined\hypertarget\hypertarget{sourcepointer-1}{}\fi
\noindent\textit{Source reference.} Original section 01ND: sources/stacks-chow/tags.txt, source lines 2154--2154. Prerequisite pointer invoked at proof line(s) 6811; it names a standard fact of a sibling chapter that the counted argument uses, and every step of the counted proof built on it is retained verbatim in the delivered body.\par\medskip
\ifdefined\hypertarget\hypertarget{sourcepointer-2}{}\fi
\noindent\textit{Source reference.} Original lemma 02L6: sources/stacks-chow/tags.txt, source lines 3397--3397. Prerequisite pointer invoked at proof line(s) 6814; it names a standard fact of a sibling chapter that the counted argument uses, and every step of the counted proof built on it is retained verbatim in the delivered body.\par\medskip
\ifdefined\hypertarget\hypertarget{sourcepointer-3}{}\fi
\noindent\textit{Source reference.} Original reference 0E3L: sources/stacks-curves/pinned-original-source.tex, source lines 6135--6135. Prerequisite pointer invoked at proof line(s) 6804; it names a standard fact of the same chapter that the counted argument uses, and every step of the counted proof built on it is retained verbatim in the delivered body.\par\medskip
\ifdefined\hypertarget\hypertarget{sourcepointer-4}{}\fi
\noindent\textit{Source reference.} Original reference 0E3J: sources/stacks-curves/pinned-original-source.tex, source lines 6032--6032. Prerequisite pointer invoked at proof line(s) 6844; it names a standard fact of the same chapter that the counted argument uses, and every step of the counted proof built on it is retained verbatim in the delivered body.\par\medskip
\ifdefined\hypertarget\hypertarget{sourcepointer-5}{}\fi
\noindent\textit{Source reference.} Original reference 0E64: sources/stacks-curves/pinned-original-source.tex, source lines 6317--6317. Prerequisite pointer invoked at proof line(s) 6801; it names a standard fact of the same chapter that the counted argument uses, and every step of the counted proof built on it is retained verbatim in the delivered body.\par\medskip
\ifdefined\hypertarget\hypertarget{sourcepointer-6}{}\fi
\noindent\textit{Source reference.} Original reference 0E63: sources/stacks-curves/pinned-original-source.tex, source lines 5987--5987. Prerequisite pointer invoked at proof line(s) 6801; it names a standard fact of the same chapter that the counted argument uses, and every step of the counted proof built on it is retained verbatim in the delivered body.\par\medskip
\ifdefined\hypertarget\hypertarget{sourcepointer-7}{}\fi
\noindent\textit{Source reference.} Original reference 0BS6: sources/stacks-curves/pinned-original-source.tex, source lines 1218--1218. Prerequisite pointer invoked at proof line(s) 6848; it names a standard fact of the same chapter that the counted argument uses, and every step of the counted proof built on it is retained verbatim in the delivered body.\par\medskip
\ifdefined\hypertarget\hypertarget{sourcepointer-8}{}\fi
\noindent\textit{Source reference.} Original reference 0E3E: sources/stacks-curves/pinned-original-source.tex, source lines 5680--5680. Prerequisite pointer invoked at proof line(s) 6830; it names a standard fact of the same chapter that the counted argument uses, and every step of the counted proof built on it is retained verbatim in the delivered body.\par\medskip
\ifdefined\hypertarget\hypertarget{sourcepointer-9}{}\fi
\noindent\textit{Source reference.} Original lemma 0B5Y: sources/stacks-chow/tags.txt, source lines 14051--14051. Prerequisite pointer invoked at proof line(s) 6800; it names a standard fact of a sibling chapter that the counted argument uses, and every step of the counted proof built on it is retained verbatim in the delivered body.\par\medskip
\ifdefined\hypertarget\hypertarget{sourcepointer-10}{}\fi
\noindent\textit{Source reference.} Original lemma 0BUG: sources/stacks-chow/tags.txt, source lines 14874--14874. Prerequisite pointer invoked at proof line(s) 6838; it names a standard fact of a sibling chapter that the counted argument uses, and every step of the counted proof built on it is retained verbatim in the delivered body.\par\medskip
\ifdefined\hypertarget\hypertarget{sourcepointer-11}{}\fi
\noindent\textit{Source reference.} Original lemma 0E8T: sources/stacks-chow/tags.txt, source lines 17828--17828. Prerequisite pointer invoked at proof line(s) 6820; it names a standard fact of a sibling chapter that the counted argument uses, and every step of the counted proof built on it is retained verbatim in the delivered body.\par\medskip
\begin{lemma}
\label{lemma-tricanonical}
Let $k$ be a field. Let $X$ be a proper scheme over $k$ of dimension $1$ with
$H^0(X, \mathcal{O}_X) = k$ having genus $g \geq 2$. Assume the singularities
of $X$ are at-worst-nodal and $\omega_X$ is ample. Then
$\omega_X^{\otimes 3}$ is very ample and $H^1(X, \omega_X^{\otimes 3}) = 0$.
\end{lemma}

\begin{proof}
Combining Varieties, Lemma
\ref{varieties-lemma-ampleness-in-terms-of-degrees-components} and
Lemmas \ref{lemma-rational-tail-negative} and \ref{lemma-rational-bridge-zero}
we see that $X$ contains no rational tails or bridges.
Then we see that $\omega_X^{\otimes 3}$ is globally generated
by Lemma \ref{lemma-contracting-rational-tails}.
Choose a $k$-basis $s_0, \ldots, s_n$ of
$H^0(X, \omega_X^{\otimes 3})$. We get a morphism
$$
\varphi_{\omega_X^{\otimes 3}, (s_0, \ldots, s_n)} :
X \longrightarrow \mathbf{P}^n_k
$$
See Constructions, Section \ref{constructions-section-projective-space}.
The lemma asserts that this morphism is a closed immersion.
To check this we may replace $k$ by its algebraic closure, see
Descent, Lemma \ref{descent-lemma-descending-property-closed-immersion}.
Thus we may assume $k$ is algebraically closed.

\medskip\noindent
Assume $k$ is algebraically closed.
We will use Varieties, Lemma
\ref{varieties-lemma-variant-separate-points-tangent-vectors}
to prove the lemma.
Let $Z \subset X$ be a closed subscheme of degree $2$ over $Z$
with ideal sheaf $\mathcal{I} \subset \mathcal{O}_X$.
We have to show that
$$
H^0(X, \mathcal{L}) \to H^0(Z, \mathcal{L}|_Z)
$$
is surjective. Thus it suffices to show that
$H^1(X, \mathcal{I}\mathcal{L}) = 0$.
To do this we will use Lemma \ref{lemma-vanishing-on-gorenstein}.
Thus it suffices to show that
$$
3\deg(\omega_X|_Y) > -2\chi(Y, \mathcal{O}_Y) + \deg(Z \cap Y)
$$
for every reduced connected closed subscheme $Y \subset X$.
Since $k$ is algebraically closed and $Y$ connected and reduced
we have $H^0(Y, \mathcal{O}_Y) = k$ (Varieties, Lemma
\ref{varieties-lemma-proper-geometrically-reduced-global-sections}).
Hence $\chi(Y, \mathcal{O}_Y) = 1 - \dim H^1(Y, \mathcal{O}_Y)$.
Thus we have to show
$$
3\deg(\omega_X|_Y) > -2 + 2\dim H^1(Y, \mathcal{O}_Y) + \deg(Z \cap Y)
$$
which is true by Lemma \ref{lemma-no-rational-tail}
except possibly if $Y = X$ or if $\deg(\omega_X|_Y) = 0$.
Since $\omega_X$ is ample the second possibility does not
occur (see first lemma cited in this proof). Finally, if
$Y = X$ we can use Riemann-Roch (Lemma \ref{lemma-rr})
and the fact that $g \geq 2$ to see that the inquality holds.
The same argument with $Z = \emptyset$ shows that
$H^1(X, \omega_X^{\otimes 3}) = 0$.
\end{proof}

\end{document}
